\section{Opis problemu}
\label{sec:problem}

\subsection{Cel ćwiczenia}
Celem ćwiczenia jest zaprojektowanie i uruchomienie w programie
LOGO!Soft Comfort układu sekwencyjnego w języku drabinkowym (LAD),
który za pomocą \emph{jednego} przycisku zasilania steruje stanem
telefonu oraz jego ekranu. O wykonywanej akcji decyduje czas
przytrzymania przycisku i bieżący stan urządzenia, a więc układ musi
posiadać pamięć -- jest to układ sekwencyjny, a nie kombinacyjny.

\subsection{Rodzaje kliknięć}
Zdarzenia generowane przez przycisk zostały podzielone na trzy klasy
zależnie od czasu $t$ przytrzymania (tab.~\ref{tab:klikniecia}).

\begin{table}[H]
\centering
\caption{Klasyfikacja kliknięć przycisku zasilania.}
\label{tab:klikniecia}
\begin{tabular}{llcl}
\toprule
Rodzaj & Oznaczenie & Czas przytrzymania & Moment rozpoznania \\
\midrule
krótkie & K & $t < \SI{0,5}{\second}$ & w chwili puszczenia przycisku \\
średnie & Ś & $\SI{0,5}{\second} \le t < \SI{2}{\second}$ & w chwili puszczenia przycisku \\
długie  & D & $t \ge \SI{2}{\second}$ & po upływie \SI{2}{\second}, jeszcze podczas trzymania \\
\bottomrule
\end{tabular}
\end{table}

Kliknięcie długie jest rozpoznawane bez czekania na puszczenie
przycisku, podobnie jak w rzeczywistych telefonach. Puszczenie
przycisku po kliknięciu długim nie generuje już żadnego zdarzenia.

\subsection{Wymagania funkcjonalne}
\begin{enumerate}
    \item[R1.] Gdy telefon jest wyłączony, kliknięcia K i Ś nie powodują
          żadnej reakcji. Kliknięcie D uruchamia telefon i włącza ekran.
    \item[R2.] Gdy telefon jest włączony, kliknięcie K przełącza stan
          ekranu (włączony/wygaszony).
    \item[R3.] Gdy telefon jest włączony, kliknięcie Ś otwiera menu
          wyłączenia. Otwarte menu zawsze wymusza włączony ekran.
    \item[R4.] Gdy menu jest otwarte, kliknięcie K anuluje menu
          i przywraca stan ekranu sprzed jego otwarcia.
    \item[R5.] Gdy menu jest otwarte, kliknięcie Ś pozostawia menu
          otwarte.
    \item[R6.] Kliknięcie D w dowolnym stanie z włączonym telefonem
          wyłącza telefon, zamyka menu i gasi ekran.
    \item[R7.] Ekran nigdy nie świeci przy wyłączonym telefonie.
    \item[R8.] Jedno naciśnięcie przycisku wywołuje dokładnie jedno
          zdarzenie (K, Ś albo D). Dłuższe trzymanie przycisku nie
          powoduje powtórnych przełączeń.
\end{enumerate}

\subsection{Wielkości wejściowe i wyjściowe}
\begin{table}[H]
\centering
\caption{Wejścia i wyjścia układu.}
\label{tab:io}
\begin{tabular}{clll}
\toprule
Symbol & Typ & Znaczenie & Wartość 1 \\
\midrule
$I_1$ & wejście & przycisk zasilania (chwilowy) & przycisk wciśnięty \\
$Q_1$ & wyjście & stan ekranu & ekran włączony \\
$Q_2$ & wyjście & stan telefonu & telefon działa \\
\bottomrule
\end{tabular}
\end{table}

\subsection{Zmienne pamięci wewnętrznej}
Do zapamiętania stanu automatu oraz do rozpoznawania zdarzeń użyto
markerów i bloków funkcyjnych (tab.~\ref{tab:pamiec}).

\begin{table}[H]
\centering
\caption{Zmienne pamięci wewnętrznej.}
\label{tab:pamiec}
\begin{tabular}{cll}
\toprule
Symbol & Znaczenie & Rola \\
\midrule
$M_1$ & stan telefonu & zmienna stanu \\
$M_2$ & zapamiętany stan ekranu & zmienna stanu \\
$M_3$ & menu wyłączenia otwarte & zmienna stanu \\
$M_4$ & zdarzenie K (krótkie kliknięcie) & impuls jednocyklowy \\
$M_5$ & zdarzenie Ś (średnie kliknięcie) & impuls jednocyklowy \\
$M_6$ & zdarzenie D (długie kliknięcie) & impuls jednocyklowy \\
$M_7$ & $I_1$ z poprzedniego cyklu & pamięć cyklu \\
$M_8$\footnotemark & $T_{001}$ z poprzedniego cyklu & pamięć cyklu \\
$M_9$ & $T_{002}$ z poprzedniego cyklu & pamięć cyklu \\
$T_{001}$ & timer TON \SI{0,5}{\second} & pomiar czasu \\
$T_{002}$ & timer TON \SI{2}{\second} & pomiar czasu \\
\bottomrule
\end{tabular}
\end{table}
\footnotetext{Na zrzutach ekranu z programu LOGO! marker ten występuje
pod oznaczeniem $M_{80}$. Pełni on rolę $M_8$ opisaną w dokumencie.}

\subsection{Założenia projektowe}
\begin{itemize}
    \item Program wykonywany jest cyklicznie przez sterownik. Przy
          przytrzymanym przycisku sygnał $I_1$ pozostaje w stanie
          wysokim przez wiele cykli, dlatego każde zdarzenie musi być
          sprowadzone do impulsu trwającego \emph{dokładnie jeden
          cykl}.
    \item Wszystkie przerzutniki RS są bezpamięciowe po restarcie
          programu (\texttt{Rem=off}), więc układ zawsze startuje
          w stanie $S_0$ (telefon wyłączony).
    \item Przy każdym przerzutniku warunki Set i Reset zostały
          dobrane tak, aby nigdy nie były aktywne jednocześnie.
          Poprawność nie zależy zatem od priorytetu wejść bloku RS.
\end{itemize}
